\scanpage{27}
continu

\[
f^\bullet:\underline{\mathrm{Et}}_S\qquad\underline{\mathrm{Et}}_T\quad .
\]

On suppose donné une donnée de ramification $\underline{R}_S=(\underline{C}_S,\underline{s}_S)$ sur $S$, une autre $\underline{R}_T=(\underline{C}_T,\underline{s}_T)$ sur $T$, et un foncteur \inserted{noté} $M\quad M\mathbin{\mathrm{x}}_S T:\underline{C}_{\unsure{S|T}}\quad\underline{C}_T$, transformant morphismes cartésiens en morphismes carétsiens et rendant commutatif de diagramme

\diagram{\begin{tikzcd}[ampersand replacement=\&,column sep=large]
\underline{C}_S \& \& \underline{C}_T\\
\underline{\mathrm{Et}}_S \& {f^\bullet} \& \underline{\mathrm{Et}}_T\quad .
\end{tikzcd}}

Considérons de plus le carré de foncteurs

\diagram{\begin{tikzcd}[ampersand replacement=\&,column sep=small,row sep=small]
\underline{C}_S \& \& \underline{C}_T\\
\underline{s}_S \& \& \underline{s}_T\\
(\mathrm{Sch})_{/S} \& \& (\mathrm{Sch})_{/T}\\
\& {X\quad X\mathbin{\mathrm{x}}_S T} \&
\end{tikzcd}}

et supposons donné un morphisme $u$ entre les deux composés:

\[
u(M):\quad\underline{s}_T(M\mathbin{\mathrm{x}}_S T)\qquad
\underline{s}_{\deleted{\unsure{T}}S}(M)\mathbin{\mathrm{x}}_S T
\]

qui, pour tout $M$ fixé\deleted{\unsure{, soit un}} sur $S'\quad\mathrm{Ob}\,\underline{\mathrm{Et}}_S$, soit un ${\unsure{F|f}}^{\bullet}(S')$-morphisme. Sous ces conditions, on définit de façon évidente un foncteur, noté $f^\bullet$ \inserted{ou $M\quad M\mathbin{\mathrm{x}}_S T$} si une confusion n'est pas à craindre, de $\mathrm{Rev}^{\underline R}(S)$ dans $\mathrm{Rev}^{\underline R_T}(T)$ et un homomorphisme de foncteurs composés

\diagram{\begin{tikzcd}[ampersand replacement=\&,column sep=small]
\mathrm{Rev}^{\underline R(S)}(S) \& \& \mathrm{Rev}^{\underline R(T)}(T)\\
\mathrm{Rev}(S) \& \& \mathrm{Rev}(T)\quad .\\[-1ex]
\& {X\quad X\mathbin{\mathrm{x}}_S T} \&
\end{tikzcd}}

Cet homomorphisme est un isomorphisme dans certains cas, par exemple sous les conditions de l'exemple 3, lorsqu'on suppose (sous les notations d'icelui) que pour tout sous-groupe $H$ de $G$, le passage au quotient $\underline{s}(M)/H$ commute au changement de base $T\quad S$; cette condition sera vérifié
% archive sheet number: 26
